В ходе работы были спроектированы четыре модели хранения данных
системы видеонаблюдения: PostgreSQL JSONB, нормализованная PostgreSQL,
вложенная MongoDB и нормализованная MongoDB. Все модели описывают одни и те
же данные: территории, зоны, камеры, аналитические события и телеметрию.

Для сравнения этих моделей был реализован стенд нагрузочного тестирования.
Для запуска прогонов и просмотра метрик использовалось веб-приложение.
Сравнение моделей выполнялось по времени записи, времени чтения и пропускной
способности.

Лучший результат по сумме метрик показала вложенная MongoDB: она дала
наибольшую пропускную способность и особенно хорошо показала себя при чтении.
PostgreSQL JSONB не обогнала вложенную MongoDB, но была немного быстрее
нормализованной PostgreSQL как на запись, так и на чтение. Нормализованная
MongoDB уменьшила дублирование данных, однако значительно проиграла всем
остальным моделям по производительности.

Таким образом, в проведенных экспериментах выиграла вложенная документная
модель MongoDB. Главная причина -- денормализация данных, которые обычно
читаются вместе в сценариях видеонаблюдения. При аналитических запросах
системе не нужно соединять несколько сущностей, а при записи сложного события
не нужно раскладывать его на несколько таблиц. За счет этого уменьшаются
накладные расходы, растет пропускная способность и снижаются задержки.
